\documentclass[a4paper, 10pt]{article}
\usepackage{scrextend}
\changefontsizes{9pt}
\input{../course-config}
\usepackage{../vendor/template/CEDT-Assignment-style}

\begin{document}
\subject
\asgntitle{X}{}{Week x - x}{\StudentID\ \StudentName}{\DefaultCollaborators}

% ================================================================================ %
\section{PUT THE SECTION HERE}
% ================================================================================ %

% ================================================================================ %
%                                    Problem 0x                                    %
% ================================================================================ %
\begin{problem}
Here's problem statement.
\begin{codingbox}
    print("You can insert code like this")
\end{codingbox}
You can insert sub-problem
\begin{subproblems}
	\item Here's sub-problem a.
	\item Here's sub-problem b.
	\item Here's sub-problem c.
\end{subproblems}
\end{problem}

\begin{solution}
	The solution goes here.
\end{solution}
% ================================================================================ %

% ================================================================================ %
%                                    Problem 0x                                    %
% ================================================================================ %
\begin{problem}
Here's the problem statement for \textbf{MANDATORY} problem.
\end{problem}

\begin{tosubmit}
	Here's the solution for the \textbf{MANDATORY} problem. \\
	ABC.
\end{tosubmit}
% ================================================================================ %

% ================================================================================ %
%                                    Problem 0x                                    %
% ================================================================================ %
\begin{problem}
Here's the problem statement for \textbf{PROVING} problem.
\end{problem}

\begin{proofbox}
	Here's the proof for the \textbf{PROVING} problem.
\end{proofbox}
% ================================================================================ %

% ================================================================================ %
%                                    Problem 0x                                    %
% ================================================================================ %
\begin{problem}
Here's a \textbf{multiple choice} problem statement. Use \verb|\choices| to lay
out the options; it defaults to 4 choices and picks 4, 2, or 1 columns on its
own depending on how long the choices are.
\begin{subproblems}
	\item Short choices sit four to a row.
	      \choices{True, False, Maybe, Unknown}
	\item Longer choices wrap down to two columns automatically.
	      \choices{The probability increases, The probability decreases, The probability stays the same, Not enough information to tell}
	\item Pass the choice count explicitly when it isn't 4.
	      \choices[6]{A, B, C, D, E, F}
\end{subproblems}
\end{problem}
% ================================================================================ %

\end{document}